% 0_abstract.tex
% LQRL: Layered Q-Value with Learning Reward for Non-Parametric Skill Evolution

We propose \textbf{LQRL} (Layered Q-Value with Learning Reward), a novel skill Q-value update rule for non-parametric skill evolution systems. In these systems, skills improve their content through experience, but standard Q-learning only rewards task success, creating a contradiction: skills that learn from failure are penalized with lower retrieval probability even though their content has improved. LQRL resolves this by decomposing the reward signal into a task component ($r_{\text{task}}$) and a learning component ($r_{\text{learning}}$). The learning reward is evaluated by an LLM judge that assesses whether skill content improved after a task, allowing skills to be rewarded for learning even when they fail. Our update rule is $Q_{\text{new}} = Q_{\text{old}} + \alpha \cdot [(1-\beta) \cdot r_{\text{task}} + \beta \cdot r_{\text{learning}}]$, where $\beta$ controls the weight of learning reward. Experiments on BigCodeBench show that $\beta=0.3$ improves success rate from 40\% to 66.7\% compared to standard Q-learning ($\beta=0$), demonstrating that decomposing rewards enables skills to learn effectively from both success and failure.
